% s-099-p099-m002; original PDF 99, 100, 101, 102, 103, 104
\Aphorism{s-099-p099-m002}{100}
\PoemSubtitle{Zwischen zwei Schulen.}
\Sperr{Schüler.} Ich soll also nicht länger dein Schüler sein?
\par

\Sperr{Meister.} Nein, das Leben wartet auf dich.
\par

\Sperr{Schüler.} Kann ich nicht leben und zugleich —
\par

\Sperr{Meister.} Ich, oder das Leben. Diese beiden Lehrer vertragen sich nicht an Einem Schüler.
\par

\Sperr{Schüler.} Lehrst du denn gegen das Leben?
\par

\Sperr{Meister.} Wir lehren dasselbe, mit entgegengesetzten Mitteln.
\par

\Sperr{Schüler.} So verzichte ich auf das Leben und bleibe bei dir.
\par

\Sperr{Meister.} Aber das Leben verzichtet nicht auf dich, und ich behalte dich nicht. Du bist der erste, den ich fortjagen muss; alle anderen liefen mir zu früh aus der Schule.
\par

\Sperr{Schüler.} Aber auch ich kann dich noch nicht entbehren.
\par

\Sperr{Meister.} Damit du es lernst, schicke ich dich fort. Auch du wirst mir fehlen, darum eben müssen wir uns trennen. Zögern wir, so begiebt sich die Komödie.
\par

\Sperr{Schüler}. Was nennst du die Komödie?
\par

\Sperr{Meister}. Dass ein Mensch nicht ohne einen anderen Menschen leben kann.
\par

\Sperr{Schüler}. Das ist die Liebe!
\par

\Sperr{Meister}. Die Liebe ist Komödie. Du verstehst, die unüberwindliche Liebe, die Liebe als Fatum und Ewigkeit. Darüber lachen wir.
\par

\Sperr{Schüler}. Und über welche Liebe darf nicht gelacht werden?
\par

\Sperr{Meister}. Darf nicht? Freund, bestehe deine Schlussprüfung besser. Auch über die erlaubte Liebe ist Lachen erlaubt, über die Liebe als Pensum und Zwischenzustand.
\par

\Sperr{Schüler}. So erlaube mir, dich noch eine Weile zu lieben.
\par

\Sperr{Meister}. Nicht eine Minute länger. Wir sind in Gefahr, uns unentbehrlich zu werden.
\par

\Sperr{Schüler}. Aber fern von dir werde ich dich um so mehr lieben.
\par

\Sperr{Meister}. So musst du schnell noch über mich lachen lernen.
\par

\Sperr{Schüler}. Das lerne ich nie: ich verehre dich.
\par

\Sperr{Meister}. Um so schlimmer; so muss ich dir auch noch verächtlich werden. — Dich aber liebe ich nur, ich verehre dich nicht. Also genügt es, dass ich über dich lache.
\par

\Sperr{Schüler}. Das verdiene ich redlich.
\par

\Sperr{Meister}. Keineswegs. Du musst noch etwas dazu thun. Verheirathe dich!
\par

\Sperr{Schüler}. Meister, ist das dein Ernst?
\par

\Sperr{Meister}. Nicht mehr noch minder als alle meine Wahrheiten. Glaube, dass du ein Weib nehmen musst, wie du an meine Philosophie glaubst. \Sperr{Schüler}. Das eine fällt schwerer als das andere.
\par

\Sperr{Meister}. Wie? fiel dir meine Philosophie zu glauben leicht? So streiche sie durch, sie ist unwahr.
\par

\Sperr{Schüler}. Dann gieb mir eine andere.
\par

\Sperr{Meister}. Das Gegentheil! Also: im Anfang war das Nichts, und das Nichts hatte Selbsterhaltungstrieb.
\par

\Sperr{Schüler}. Folglich wurde es zu Etwas.
\par

\Sperr{Meister}. Gut! du nimmst dir schon gewisse Freiheiten von der Logik. Aber du hast noch nicht Freiheit genug, also verheirathe dich!
\par

\Sperr{Schüler}. Ich hörte sagen, Ehe sei Selbstbindung.
\par

\Sperr{Meister}. Ich glaube gar, du willst die Freiheit im Freien lernen.
\par

\Sperr{Schüler. Hab’ ich doch denken von einem Denker} gelernt!
\par

\Sperr{Meister}. Die schlechteste Schule! Im Leben ist die kürzeste Linie die krumme. Bei mir bekamst du Wahrheiten oder wenigstens Zusammenhänge, nun sollst du dich auf Fragmente und Ironie verstehen lernen. Oh wie wenig bist du auf das Leben vorbereitet!
\par

\Sperr{Schüler}. Ich bin dir dafür dankbar!
\par

\Sperr{Meister}. Und wenn der erste Fehlgriff kommt, wirst du dann noch dankbar sein? Du, der den Fehlgriff für Dasjenige hält, was man vermeiden soll! Denn du willst ja nicht einmal heirathen!
\par

\Sperr{Schüler}. So werde ich es thun, dir zu lieb.
\par

\Sperr{Meister}. Nicht mir zu lieb, dir zu leid! damit ich zu lachen habe, wenn du fehlgreifst!
\par

\Sperr{Schüler}. Bist du denn schadenfroh?
\par

\Sperr{Meister}. Mich freut es, wenn Einer sein Unerlässliches erfährt. Glück in der Ehe wünschen dir schon Allzuviele.
\par

\Sperr{Schüler}. Ist Glück verboten? \Sperr{Meister}. Nicht \Sperr{Jedem,} aber \Sperr{dir!} Ich sehe, du hast dein Abc verlernt! Glück bindet; in ihm bleibt man liegen und geht seinem Ziele verloren. Alle Abtrünnigen sind glücklich.
\par

\Sperr{Schüler.} Bin ich ein Gesendeter?
\par

\Sperr{Meister.} Das soll sich zeigen. Willst du deine Sendung in der Tasche haben, um auf Glück zu verzichten? Dass ich dich vom Glücke wegrufe, ist deine erste Berufung; schweigend erwarte die zweite!
\par

\Sperr{Schüler.} Ich will mich enthalten.
\par

\Sperr{Meister.} Nein! du sollst Misslingen und Enttäuschung aufsuchen.
\par

\Sperr{Schüler}. Und darin mein Glück finden.
\par

\Sperr{Meister}. Du meinst, in den Folgerungen, zu denen Missgeschick und Enttäuschung zu berechtigen scheinen? Sobald du innewirst, dass es dich erleichtert, das Leben zu lästern, ist es Zeit, ein mässiges Glück aufzusuchen, — nur soviel, dir die Bitterkeit aus dem Herzen herauszulächeln.
\par

\Sperr{Schüler}. Nicht mehr? nie mehr?
\par

\Sperr{Meister}. Eines Tages auch mehr. Sobald du dir bewiesen hast, Abschied nehmen zu können, darfst du auch bleiben. Sobald das Glück dich nicht mehr zu sich hinabzieht, deine Erkenntnisse nicht mehr fälscht, sobald du das Äusserste gegen dich verhängen kannst, ohne mit der Wimper zu zucken — bist du frei und Meister.
\par

\Sperr{Schüler}. Muss ich also unpersönlich werden?
\par

\Sperr{Meister}. Das Persönliche, das dich hemmt, trübt, zum Parteigänger deiner selbst macht, musst du überwinden. Aber ist es überwunden, so hast du eine neue Persönlichkeit, der das Nothwendige leicht fällt.
\par

\Sperr{Schüler}. Aber mit dem Allem hat doch die Ehe nichts zu schaffen? \Sperr{Meister}. Kennst du einen engeren Raum, wo alles beisammen ist: Fehlgriff, Enttäuschung über das ganze Leben, Heilung durch Sonnenblicke des Glücks und Erziehung zur Unabhängigkeit?
\par

\Sperr{Schüler}. Wenn ich nun selbst in die Ehe trete, bin ich denn auch darin?
\par

\Sperr{Meister}. Du musst ein Weib nehmen, das dich überlebt und bei Lebzeiten nicht loslässt.
\par

\Sperr{Schüler}. Und wann soll ich philosophiren?
\par

\Sperr{Meister}. In geraubten Augenblicken. Lerne in Blitzen und Hammerschlägen denken. Was bisher Philosophie hiess, ist ein Kirchhof von Gedanken, die schon todt waren, als man sie zu Papier brachte. Gedanken sind eine Secunde lang wahr; in dieser muss man sie fassen. Um Philosoph zu werden, heirathe.
\par

\Sperr{Schüler}. Um ein Weib zu nehmen, muss ich es lieben.
\par

\Sperr{Meister}. Das sollst du, sonst ist ja kein Fehlgriff wahrscheinlich. Du sollst an der fixen Idee leiden: das ist sie, die Einzige, ohne sie kann ich nicht leben — denn das bedeutet, dass sie die Einzige ist, mit der du nicht leben kannst.
\par

\Sperr{Schüler}. Also die Komödie, die sich als Tragödie entpuppt.
\par

\Sperr{Meister}. Grosse Worte für kleine Dinge. Tragödie? wenn der Held nur ein Mensch ist?
\par

\Sperr{Schüler}. Und der Vernichter nur \Sperr{ein Weib!} — Ich bin bereit zu fallen, aber nicht in diesem Kampfe, der eigentlich eine langsame Abnutzung ist, nicht gegen diesen Gegner.
\par

\Sperr{Meister}. Bist du wählerisch in den Anlässen, dich zu opfern? Hier bietet sich ein Untergang, und du fragst noch, ob er dir gilt? Du willst dich nicht umsonst geben; kennst du aber deinen Preis? — Oh Mensch! du kennst kaum das Leben, und hängst schon daran! — Du suchst \Sperr{deinen} Feind, nicht den ersten besten. Kann er sich dir deutlicher legitimiren, als dadurch, dass er ein Weib ist, das du liebst?
\par

\Sperr{Schüler}. Du sprichst über die Liebe wie ein Wissender.
\par

\Sperr{Meister}. Du wie ein Liebender. Du liebst ein Weib, also hast du mich längst verlassen. Gehe nun von mir zu ihm.
\par

\Sperr{Schüler}. Lass mich bei dir bleiben! Du schützest mich vor diesem Weibe!
\par

\Sperr{Meister}. Wer hier dein Schutz sein wollte, an dem nähme deine „Sendung“ Rache! Und ich fürchte diese eine Art Rache.
\par

\Sperr{Schüler}. So wärst du auf diese eine Art — feige?
\par

\Sperr{Meister}. Dafür verachtest du mich. Vielleicht nur einen Augenblick, aber der hilft uns Beiden. Leb’ wohl — und gehorche!
\par

\Sperr{Schüler}. Leb’ wohl.
\par

